  \chapter{Analýza právnych požiadaviek na bezpečnosť ISVS}

% src: https://tex.stackexchange.com/a/10012
\newcommand\lawsubtitle[1]{\begin{flushleft}\small\textit{#1}\end{flushleft}}

\section{Zákon č. 69/2018 Z. z.}
\lawsubtitle{Zákon o kybernetickej bezpečnosti a o zmene a doplnení niektorých zákonov}

Zakladným zákonom (\textit{lex generalis}), ktorý definuje riešenie kybernetickej (a informačnej) bezpečnosti v SR, je
Zákon č. 69/2018 Z. z. (\acrshort{ZoKB}). Bol prijatý v roku 2018 a prešiel niekoľkými novelizáciami \textit{373/2018 Z. z.},
\textit{134/2020 Z. z.}, \textit{287/2021 Z. z.}, \textit{55/2022 Z. z.} a \textit{231/2022 Z. z.} Implementuje európsku
smernicu \acrshort{NIS}, a preto explicitne hovorí o bezpečnosti informačných systémov a sieti vo všeobecnosti, v porovnaní
s \acrshort{ZoITVS}, ktorý hovorí o bezpečnosti informácie v informačných technologiách verejnej správy.

\subsection{Pôsobnosť zákona vo vzťahu k mestám a obciam, prípadne ďalším organizáciám}
V \acrshort{ZoKB} sú explicitne vymenované povinnosti \acrfull{NBU}(ďalej Úrad) a ďalších štátnych inštitúcií (ďalej orgány
verejnej moci). Podobne sú explicitne vymenované povinnosti poskytovateľov digitálnych služieb a prevádzkovateľov základných
služieb. Základnú službu a prevádzkovateľa definuje \acrshort{ZoKB} v \S 3 písm. l) a m):
\begin{quotation}
\begin{enumerate}[a)] \setcounter{enumi}{11}
  \item základnou službou služba, ktorá je zaradená v zozname základných služieb a
  \begin{enumerate}[1.]
    \item závisí od sietí a informačných systémov a je činnosťou aspoň v jednom sektore alebo podsektore podľa prílohy č. 1, alebo
    \item je prvkom kritickej infraštruktúry, \textsuperscript{9)}
  \end{enumerate}

  \item prevádzkovateľom základnej služby orgán verejnej moci alebo osoba, ktorá prevádzkuje aspoň jednu službu podľa písmena l),
\end{enumerate}
\end{quotation}
V prílohe č. 1 definuje sektory (a k ním príslušné podsektory) ako bankovníctvo, doprava, energetika, verejná správa a ďalšie.
Prevádzkovateľmi základnej služby sú orgány verejnej moci, mesta, obce a ďalšie organizácie. 

Zoznam prevádzkovateľov základných služieb je udržiavaný a spravovaný \acrshort{NBU} a spôsob zaradenia do zoznamu je obsiahnutý
\acrshort{ZoKB} v \S 3, ods (1) písm. k):
\begin{quotation}
\begin{enumerate}[a)] \setcounter{enumi}{10}
  \item na základe oznámenia ústredného orgánu, prevádzkovateľa základnej služby, poskytovateľa digitálnej služby alebo z
  vlastnej iniciatívy určuje
  \begin{enumerate}[1.]
    \item základnú službu a zaraďuje ju do zoznamu základných služieb,
    \item digitálnu službu a zaraďuje ju do zoznamu digitálnych služieb,
    \item poskytovateľa digitálnej služby a zaraďuje ho do registra poskytovateľov digitálnych služieb,
    \item prevádzkovateľa základnej služby a zaraďuje ho do registra prevádzkovateľov základných služieb  
  \end{enumerate}
\end{enumerate}
\end{quotation}
Pochybnosti, či je organizácia prevádzkovateľom základnej služby, vyriešilo \acrshort{NBU} vlastnou iniciatívou, keď zaradilo
veľkú časť miest a obcí na zoznam prevádzkovateľov základných služieb.

\subsection{Povinnosti prevádzkovateľa základnej služby}
Povinnosti prevádzkovateľa základnej služby sú špecifikované \acrshort{ZoKB} v \S 19. Prevádzkovateľ základnej služby je
povinný prijať a dodržiavať 12 mesiacov odo dňa oznámenia o zaradení do registra prevádzkovateľov základných služieb, všeobecné
bezpečnostné opatrenia (\acrshort{ZoKB} a jeho vykonávacími predpismi) a sektorové bezpečnostné opatrenia (\acrshort{ZoITVS}
a jeho vykonávacími predpismi). Úroveň bezpečnostných opatrení musí zohľadňovať klasifikáciu a kategorizáciu systémov\footnote{Úroveň
vyplýva zo závažností rizík vyplývajúcich z hrozieb voči systémom a aktívam organizácie.}. Prijatie bezpečnostných opatrení
závisí od výsledkov analýzy rizík\footnote{\acrshort{ZoKB} implicitne predpokladá o už vykonanej analýze rizík.}.

\acrshort{ZoKB} rámcovo stanovuje všeobecné bezpečnosné opatrenia, ktoré sú závezné pre všetky informačné systémy a siete.
Tieto bezpečnostné opatrenia majú pokrývať všetky oblasti \acrshort{KIB}. \acrshort{ZoKB} explicitne vymenúva tieto oblasti
a minimálne bezpečnostné opatrenia v \S 20:
\begin{quotation}
\begin{enumerate}[(1)] \setcounter{enumi}{2}
  \item Bezpečnostné opatrenia sa prijímajú a realizujú najmä pre oblasť
  \begin{enumerate}[a)]
  \item organizácie kybernetickej bezpečnosti a informačnej bezpečnosti,
    \item riadenia rizík kybernetickej bezpečnosti a informačnej bezpečnosti,
    \item personálnej bezpečnosti,
    \item riadenia prístupov,
    \item riadenia kybernetickej bezpečnosti a informačnej bezpečnosti vo vzťahoch s tretími stranami,
    \item bezpečnosti pri prevádzke informačných systémov a sietí,
    \item hodnotenia zraniteľností a bezpečnostných aktualizácií,
    \item ochrany proti škodlivému kódu,
    \item sieťovej a komunikačnej bezpečnosti,
    \item akvizície, vývoja a údržby informačných sietí a informačných systémov,
    \item zaznamenávania udalostí a monitorovania,
    \item fyzickej bezpečnosti a bezpečnosti prostredia,
    \item riešenia kybernetických bezpečnostných incidentov,
    \item kryptografických opatrení,
    \item kontinuity prevádzky,
    \item auditu, riadenia súladu a kontrolných činností.
  \end{enumerate}

  \item Bezpečnostné opatrenia musia zahŕňať najmenej
  \begin{enumerate}[a)]
    \item určenie manažéra kybernetickej bezpečnosti, ktorý je pri návrhu, prijímaní a presadzovaní bezpečnostných opatrení
    nezávislý od štruktúry riadenia prevádzky a vývoja služieb informačných technológií a ktorý spĺňa znalostné štandardy
    pre výkon roly manažéra kybernetickej bezpečnosti,
    \item detekciu kybernetických bezpečnostných incidentov,
    \item evidenciu kybernetických bezpečnostných incidentov,
    \item postupy riešenia a riešenie kybernetických bezpečnostných incidentov,
    \item určenie kontaktnej osoby pre prijímanie a evidenciu hlásení,
    \item pripojenie do komunikačného systému pre hlásenie a riešenie kybernetických bezpečnostných incidentov a centrálneho
    systému včasného varovania.
  \end{enumerate}
\end{enumerate}
\end{quotation}
Podrobná špecifikácia bezpečnostných opatrení, kategórií a ďalších povinností prevádzkovateľa základnej služby je uvedená
vo Vyhláška Národného bezpečnostného úradu č. 362/2018 Z.z.

Riešenie bezpečnostných incidentov ukladá \acrshort{ZoKB} v \S 24. Prevádzkovateľ základnej služby je povinný nahlasovať
závažné bezpečnostné incidenty. Ďalšie povinnosti sú stanovené v \S 27 je popísane, kde \acrshort{NBU} vydáva výstrahy a
varovania pred hrozbami a ukladá prevádzkovateľom základnej služby riešiť incident alebo vykonať reaktívne (preemptívne)
opatrenia. Zo \acrshort{ZoITVS} vyplýva, že prevádzkovateľ by mal mať definovaný postup riešenia bezpečnostných incidentov,
vrátane nahlasovania, vyhodnocovania a prijímania adekvátnych bezpečnostných opatrení.

\acrshort{ZoKB} síce v \S 24 stanovuje povinnosť prevádzkovateľovi základnej služby do dvoch rokoch po zaradení do zoznamu
sa podrobiť auditu. Avšak pre kategóriu I. Vyhláška Národného bezpečnostného úradu č. 362/2018 Z.z (doplnená a novelizovaná
Vyhláškou Národného bezpečnostného úradu č. 264/2023 Z.z.) túto povinnosť neustanovuje\footnote{Je iba odporúčaním.}.

\acrshort{ZoKB} ukladá v \S 19 ukladá zmluvné povinnosti prevádzkovateľa základnej služby voči tretím stranám: 
\begin{quotation}
\begin{enumerate}[(1)] \setcounter{enumi}{1}
  \item Prevádzkovateľ základnej služby je povinný pri výkone činnosti, ktorá priamo súvisí s dostupnosťou, dôvernosťou
  a integritou prevádzky sietí a informačných systémov prevádzkovateľa základnej služby prostredníctvom tretej strany,
  uzatvoriť zmluvu o zabezpečení plnenia bezpečnostných opatrení a notifikačných povinností podľa tohto zákona počas celej
  doby výkonu tejto činnosti; pri uzatvorení zmluvy sa vykonáva analýza rizík.
  \end{enumerate}
\end{quotation}
Podľa nasledujúceho odseku (3) nie je potreba uzavrieť takúto zmluvu, ak je tretia strana prevádzkovateľom základnej alebo
digitálnej služby, alebo ak je riziko pri vykonávaní činnosti voči tretej strane nízke. Ak nie je možné uzavrieť zmluvu
podľa odseku (2), tak je prevádzkovateľ základnej služby povinný informovať tretiu stranu v nevyhnutnom rozsahu o bezpečnostnom
incidente.

Nakoniec, podľa \acrshort{ZoKB} môže ústredný orgán definovať špecifické bezpečnostné opatrenia a vyžadovať ich zavedenie
v sektore/podsektore, za ktorý zodpovedá. Za sektor verejnej správy a podsektor informačných systémov verejnej správy zodpovedá
\acrfull{MIRRI}\footnote{V samotnej prílohe č. 1 \acrshort{ZoKB} je uvedený Úradu podpredsedu vlády pre investície a informatizáciu, \acrshort{MIRRI}
je jeho nástupníckym orgánom.}.
